% LSPV-4 REPORT::TWOSTEPTRACE1::2026-10-14T19:33:42+09:00
% COMPACT GITHUB PUBLICATION EDITION. The complete 6-page local LaTeX manuscript, source code, raw experiments and original idea register are retained separately.
% PRELITERATURE IDEAS SHA256: 09741a8ed5a5a582f2e99458d4a4b60d46cd3f0071178358ba407b6c5bb3a5bf
% CODE SHA256: 3fb04d6288bc80850522edc500de8e8a1ee845af52a72b1d242d03990ec15f63
\documentclass{article}
\usepackage{amsmath,amssymb,amsthm}
\newtheorem{theorem}{Theorem}
\title{TWOSTEPTRACE1: Adaptive Two-Read Memory Extensions and Degree Nonheredity}
\author{LSPV-4}
\date{14 October 2026}
\begin{document}\maketitle
\noindent\textbf{Status: PROVED/EXACT for the restricted explicit source, pending independent review. P versus NP remains OPEN.} This is a compact independent publication edition; the complete source and raw data remain available in the exact local research bundle. No arbitrary-algorithm separation or fast full Cook--Mertz simulation is asserted.

\section{CORPUS REVIEWED}
Live GitHub governance, AGENTS/CODEX, repository research skill, truncated root main.tex, main4 ARC, main5 WRAPPER, original SOURCELDEGATE2, COMPLETIONCUT1, LIVEPREFIX1, LSPV-3 FIBERCORE1/HALLMASK1, LSPV-5 FRONTIERRANK1 and the original Drive SIM-4-0008 were read. Local original GROUPEDREAD1/CANONCOUNT1 manuscripts and LSPV-1 ACCESSMARGIN1 were also inspected. The 2026 explicit-tree catalytic result was already in SOURCELDEGATE2 and is not claimed new.

\section{IDEAS PROPOSED BEFORE LITERATURE}
The original preregistered CANDIDATES, persisted before external literature and not retroactively edited, were: (A) source-defined two adaptive reads and writes using trace-polynomial bit extraction (missing off-grid degree budget; falsifier exhaustive packed replay); (B) delayed-write chronological trace register (missing polynomial acquisition and finite iterability; falsifier repeated accesses over b steps); (C) compare packed field interpolation with canonical off-grid lookups (missing full cost/memory control; falsifier compact-memory or off-grid controls); (D) completion-time bitstream wrapper (missing a compulsory advance-commitment obligation; falsifier post-release search). After literature A survived for two steps, B and C remain partial, D does not establish SAT hardness.

\section{LITERATURE CHECK AND REVISIONS}
Williams (STOC 2025), Theorems 1.1 and 3.1, proves square-root-space simulation for multitape TMs, without preserving the time bound necessary for our FLMV endpoint. Remark 1.6 excludes unrestricted random access. Theorem 5.1 gives a length-restricted matched-oracle version; the full simulator is described as nonrelativizing. Section 5 identifies Tree Evaluation and off-grid evaluation bottlenecks. Cook--Mertz (STOC 2024) provides the tree algorithm. Henzinger--Pyne--Ragavan (Feb. 2026) and Asadi--Cleve (Apr. 2026) provide catalytic explicit-tree tradeoffs; none transfers automatically to polynomial time in a succinctly unfolded computation DAG. No theorem supplies the required single model-matched near-time low-space realization.

\section{CONSTRUCTION OR PROOF ATTEMPT}
Fix h>=2, b=2^h, S=GF(2^h) contained in K=GF(2^{4h}); hence |K|=b^4. Let e_i be an F_2-basis of S, with trace-dual d_i. Pack b Boolean cells into t=ceil(b/h) words U_q in S, padding unused coordinates to zero. Define
\[
E_i(Z)=\sum_{j=0}^{h-1}(d_i Z)^{2^j},\qquad
\chi_k(A)=\prod_{i=0}^{h-1}\bigl(k_iE_i(A)+(1-k_i)(1+E_i(A))\bigr),
\]
\[
R(A,U)=\sum_{k<b}\chi_k(A)E_{k\bmod h}(U_{\lfloor k/h\rfloor}),
\]
\[
W_q(A,U,v)=\sum_{i<h}e_i\{E_i(U_q)+[qh+i<b]\chi_{qh+i}(A)(v+E_i(U_q))\}.
\]
All additions are characteristic two; k_i is the i-th binary digit of k. On legal S-grid inputs E_i extracts the i-th Boolean coordinate and exactly one chi_k equals one.

\begin{theorem}[Two source-dependent transitions]
For every legal A,U in S, the following arithmetic circuit agrees with a Boolean machine that first reads tape[A], flips that cell, changes the low address bit iff the first read was one, then reads and flips tape at the modified address:
\[
r_0=R(A,U),\quad U'_q=W_q(A,U,1+r_0),\quad
A'=A+e_0r_0,\quad r_1=R(A',U'),\quad
U''_q=W_q(A',U',1+r_1).
\]
The formal total degree of every output is at most
\[
D_2=(2h^2+4h+1)2^{2h-2}<|K|-1=b^4-1.
\]
For b>=8, D_2<b^3-1. Given a certified field/subfield representation and trace basis, the circuit has O(bh^2+h^3) arithmetic gates and O(b+poly(h)) live bits for sequential evaluation.
\end{theorem}
\begin{proof}
Grid correctness follows from the trace-dual identities, exclusive selectors, and two elementary read/flip operations. Set e=2^{h-1}. Then deg(E_i)<=e, deg(chi_k)<=he, deg(r_0)<=(h+1)e, deg(U'_q)<=(2h+1)e, and deg(A')<=(h+1)e. Substitution bounds second selectors by h(h+1)e^2 and extracted tape bits by (2h+1)e^2. Hence deg(r_1)<=(h^2+3h+1)e^2 and deg(U''_q)<=(2h^2+4h+1)e^2. At h=2 the latter is 68<255; for h>=3 compare 2h^2+4h+1 with 4*2^h. The circuit performs O(bh^2) source arithmetic after separately charged field preprocessing; no claim is made that brute-force field construction meets a near-time contract. Live input, intermediate and output tape words total O(b) bits.
\end{proof}

\begin{theorem}[Exact degree nonheredity]
For b=2^h>=4 and S=GF(b) inside K=GF(b^2), P(Z)=Z^{b+1} has degree b+1<|K|-1 and restricts to z->z^2 on S. Yet its literal self-composition has degree (b+1)^2>|K|-1 and agrees on S with Z^4. For b=4, it is the identity on S.
\end{theorem}
\begin{proof}
On S, z^b=z. Thus P(z)=z^2 and P(P(z))=z^4. The formal degree inequalities are immediate. Over GF(4), z^4=z.
\end{proof}
Large unreduced degree is not a lower bound: here a small normalization exists. The theorem only refutes automatic hereditary preservation of a syntactic degree promise.

\paragraph{Joint-resource obstruction.} Two explicit accesses use O(b polylog b) bit operations after field preprocessing. Literal repetition over Theta(b) steps can cost O(b^2 polylog b), before any recursive tree evaluation. A binary DAG with H+1 nodes may unfold to 2^{H+1}-1 tree nodes. A method polynomial in explicit-tree size is not necessarily polynomial in succinct H. Nothing here constructs an all-input resource contractor or a near-time, small-space one-machine simulator. The current FLMV forbidden target remains conditional on such a compiler.

\section{TESTS AND OUTCOME}
Ten independent checkpointed Python batches, seed 202610141933, completed 11,410 exact assertions and zero failures: degree 48; exhaustive b=4 320; exhaustive b=8 in four shards 10,240; off-grid 201; Horn, 2-SAT, affine, padding, DEAD/COPY/DELAY/RECOMP and wrapper controls 564; unfolded tree 6; nonheredity 31. GF(256) modulus 0x11b and GF(4096) modulus 0x1009. A nonlegal grid input with address 185 and packed values (175,212) returns field 255, underscoring that the theorem promises only legal-grid agreement. Experiments are COMPUTATIONAL; independent adversarial review remains pending.

\paragraph{Duplication, contradiction and handoff.} GROUPEDREAD1 established one step; this adds an adaptive second step. SOURCELDEGATE2 already recorded explicit-tree blowup. The root main.tex is truncated, lacks balanced LSPV-4 ownership markers and must remain untouched. LSPV-1 should charge trace compiler, field construction, all physical accesses, and complete implicit tree traversal on the same output machine. LSPV-3 should use easy source witness replay as a counterexample to inferring hardness from address dependence. LSPV-5 should not promote the local trace to universal certificate existence. Completion-time unpredictability concerns pre-release commitments; LIVEPREFIX1 shows post-release prefix search remains conditionally polynomial under P=NP. No cryptographic assumption or nonalgebrization result is used.

\paragraph{Next target.} Construct a source-checkable b-step trace compiler with degree below polynomial field cardinality and near-time off-grid evaluation of the complete implicit tree, or prove a narrow source-scoped obstruction. P versus NP remains OPEN.

\begin{thebibliography}{5}
\bibitem{w} R. Williams, Simulating Time With Square-Root Space, STOC 2025.
\bibitem{c} J. Cook and I. Mertz, Tree Evaluation Is in Space O(log n log log n), STOC 2024.
\bibitem{a} V. Asadi and R. Cleve, Polynomial-Time Almost Log-Space Tree Evaluation by Catalytic Pebbling, arXiv:2604.02606, 2026.
\bibitem{h} A. Henzinger, E. Pyne and S. Ragavan, Catalytic Tree Evaluation From Matching Vectors, arXiv:2602.14320, 2026.
\bibitem{f} L. Fortnow, R. Lipton, D. van Melkebeek, A. Viglas, Time-Space Lower Bounds for Satisfiability, J. ACM 52 (2005).
\end{thebibliography}
\end{document}
% END_OF_VERIFIED_REPORT::LSPV-4::TWOSTEPTRACE1